\section{Obiekty nad ustaloną bazą i kategorie przecinkowe}\label{sec:slice}

W geometrii często nie rozważa się po prostu przestrzeni $X$, lecz przestrzeń $X$ \emph{nad bazą} $B$, czyli odwzorowanie $p:X\to B$. Na przykład wiązka wektorowa ma projekcję na rozmaitość bazową. Naturalną kategorią dla takich danych jest kategoria obiektów nad $B$.

\begin{definicja}[Kategoria obiektów nad ustaloną bazą, slice category]\label{def:slice}
Niech $\mathcal C$ będzie kategorią i $B\in\Ob(\mathcal C)$. \emph{Kategorią obiektów nad $B$}, oznaczaną $\mathcal C/B$, nazywamy kategorię, której obiektami są morfizmy $p:X\to B$ w $\mathcal C$, a morfizm z $(X,p)$ do $(Y,q)$ jest morfizmem $u:X\to Y$ w $\mathcal C$ spełniającym $q\circ u=p$. Składanie i identyczności dziedziczymy z $\mathcal C$.
\end{definicja}

Sprawdźmy aksjomaty: $p\circ\id_X=p$, a jeśli $q\circ u=p$ i $r\circ v=q$, to
\[
r\circ(v\circ u)=(r\circ v)\circ u=q\circ u=p.
\]
Złożenie pozostaje morfizmem nad $B$. Łączność i prawa identyczności wynikają z aksjomatów $\mathcal C$. Funktor zapominający do $\mathcal C$ wysyła $(X,p)$ na $X$, a $u$ na $u$; zachowuje złożenia i identyczności oraz jest wierny, ponieważ na zbiorach morfizmów działa przez inkluzje.

\begin{stwierdzenie}[Obiekt początkowy i końcowy]\label{cat:prop:slice-initial-terminal}
Obiektem końcowym $\mathcal C/B$ jest $(B,\id_B)$. Jeżeli $\mathcal C$ ma obiekt początkowy $I$, to jedyna strzałka $i_B:I\to B$ jest obiektem początkowym $\mathcal C/B$.
\end{stwierdzenie}
\begin{proof}
Morfizm z $(X,p)$ do $(B,\id_B)$ musi spełniać $\id_B\circ u=p$, czyli $u=p$. Ta strzałka istnieje i jest jedyna. Z początkowości $I$ istnieje dokładnie jedna strzałka $i_X:I\to X$. Ponieważ $p\circ i_X$ i $i_B$ mają źródło $I$ i cel $B$, są równe. Zatem $i_X$ jest morfizmem nad $B$; każdy inny musiałby być inną strzałką $I\to X$, co jest niemożliwe.
\end{proof}

\begin{przyklad}
W $\Set/B$ obiektem początkowym jest $\varnothing\to B$, a końcowym $\id_B$. Jeśli $B=\varnothing$, istnienie funkcji $X\to B$ wymusza $X=\varnothing$.
\end{przyklad}

Obiekt kategorii $\mathcal C/B$ to więc cała strzałka, a nie samo $X$. Morfizm wyznacza przemienny diagram
\[
\begin{CD}
X @>{u}>> Y\\
@V{p}VV @VV{q}V\\
B @>>{\id_B}> B.
\end{CD}
\]

\begin{przyklad}[Zbiory nad zbiorem bazowym]
W $\Set/B$ obiektami są funkcje $p:X\to B$. Możemy rozumieć je jako \emph{rodziny zbiorów} $(X_b)_{b\in B}$, gdzie $X_b=p^{-1}(\{b\})$. Morfizm $u:(X,p)\to(Y,q)$ zachowuje włókna: z $p(x)=b$ wynika $q(u(x))=b$, więc $u$ wysyła $X_b$ do $Y_b$. Odwrotnie, dla rodziny funkcji $u_b:X_b\to Y_b$ określamy $u(x)=u_{p(x)}(x)$. Jest to funkcja dobrze określona, ponieważ $x\in X_{p(x)}$, a jej wartość należy do $Y_{p(x)}$, więc $q(u(x))=p(x)$. Ograniczenia $u$ do włókien są zadanymi $u_b$; jednoznaczność wynika z tego, że każdy $x$ należy do dokładnie jednego włókna.

Każda rodzina $(S_b)_{b\in B}$ powstaje, z dokładnością do kanonicznych bijekcji włókien, z projekcji
\[
\{(b,s):b\in B,\ s\in S_b\}\longrightarrow B,\qquad (b,s)\longmapsto b.
\]
Dla pierwotnego $p:X\to B$ odwzorowania $x\mapsto(p(x),x)$ i $(b,x)\mapsto x$ są wzajemnie odwrotne nad $B$. Są naturalne, bo dla $u$ nad $B$ zachodzi $(q(u(x)),u(x))=(p(x),u(x))$. Ograniczenie do włókien zachowuje identyczności i złożenia punktowo; konstrukcja sumy rozłącznej działa na morfizmach przez $(b,s)\mapsto(b,u_b(s))$, więc również zachowuje identyczności i złożenia. Bijekcja włókna z $S_b$ jest naturalna, gdyż $(b,s)\mapsto s$ komutuje z każdą funkcją $u_b$. Otrzymujemy równoważność kategorii $\Set/B$ i kategorii rodzin zbiorów indeksowanych przez $B$.
\end{przyklad}

\begin{stwierdzenie}[Pullback jako iloczyn nad bazą]\label{cat:prop:slice-product}
Niech $p:X\to B$ i $q:Y\to B$ będą obiektami kategorii $\mathcal C/B$. Ich iloczyn istnieje wtedy i tylko wtedy, gdy istnieje pullback $X\times_B Y$ w $\mathcal C$ (definicja~\ref{def:pullback}). Wówczas wraz ze wspólną projekcją $r=p\circ\pi_X=q\circ\pi_Y$ jest on ich iloczynem nad $B$.
\end{stwierdzenie}
\begin{proof}
Morfizmy rzutów $\pi_X:X\times_B Y\to X$ i $\pi_Y:X\times_B Y\to Y$ są morfizmami w $\mathcal C/B$, ponieważ $p\pi_X=r=q\pi_Y$. Niech $(Z,t)$ będzie obiektem nad $B$ z morfizmami $u:Z\to X$, $v:Z\to Y$ nad $B$. Wówczas $pu=t=qv$, więc z uniwersalności pullbacku otrzymujemy dokładnie jeden $h:Z\to X\times_B Y$ o własnościach $\pi_Xh=u$, $\pi_Yh=v$. Zachodzi $rh=p\pi_Xh=pu=t$, a zatem $h$ jest także morfizmem nad $B$. Każdy morfizm nad $B$ spełniający te same warunki jest równy $h$ z jednoznaczności pullbacku.

Odwrotnie, niech $(P,r)$ z rzutami $\pi_X,\pi_Y$ będzie iloczynem nad $B$. Rzuty spełniają $p\circ\pi_X=r=q\circ\pi_Y$. Dla dowolnej pary strzałek $u:Z\to X$, $v:Z\to Y$ o własności $p\circ u=q\circ v$ oznaczmy wspólną strzałkę przez $t$. Wówczas $u,v$ są morfizmami nad $B$ z $(Z,t)$, więc istnieje jedyny morfizm $h:(Z,t)\to(P,r)$ o zadanych złożeniach z rzutami. Jeśli dowolna strzałka $k:Z\to P$ w $\mathcal C$ ma te same złożenia, to $r\circ k=p\circ\pi_X\circ k=p\circ u=t$. Jest więc morfizmem nad $B$ i $k=h$. To dowodzi własności uniwersalnej pullbacku.
\end{proof}

\begin{przyklad}[Włókna iloczynu nad bazą]
Dla $p:X\to B$ i $q:Y\to B$ w $\Set$ mamy $(X\times_B Y)_b\cong X_b\times Y_b$. Para $(x,y)$ należy do włókna nad $b$ wtedy i tylko wtedy, gdy $p(x)=b=q(y)$. Ta równość wyjaśnia, dlaczego pullback nazywa się \emph{iloczynem włóknistym}.
\end{przyklad}

\begin{definicja}[Ogólna kategoria przecinkowa]\label{cat:def:comma}
Jeżeli $F:\mathcal A\to\mathcal C$ i $G:\mathcal D\to\mathcal C$ są funktorami, to \emph{kategoria przecinkowa} $(F\downarrow G)$ ma obiekty $(A,D,u)$, gdzie $A\in\Ob(\mathcal A)$, $D\in\Ob(\mathcal D)$ oraz $u:F(A)\to G(D)$. Morfizm $(a,d):(A,D,u)\to(A',D',u')$ składa się ze strzałek $a:A\to A'$ i $d:D\to D'$ spełniających $G(d)\circ u=u'\circ F(a)$. Identyczności i złożenia określamy współrzędnie.
\end{definicja}

Para $(\id_A,\id_D)$ spełnia warunek przemienności, gdyż funktory zachowują identyczności. Jeżeli dodatkowo $(a',d'):(A',D',u')\to(A'',D'',u'')$ jest morfizmem, to
\[
\begin{aligned}
G(d'\circ d)\circ u
&=G(d')\circ G(d)\circ u\\
&=G(d')\circ u'\circ F(a)\\
&=u''\circ F(a')\circ F(a)\\
&=u''\circ F(a'\circ a).
\end{aligned}
\]
Zatem złożenie jest morfizmem. Łączność i prawa identyczności zachodzą w obu współrzędnych. Rzuty na $\mathcal A$ i $\mathcal D$ są funktorami, ponieważ identyczności i złożenia określono współrzędnie.

\begin{przyklad}[Obiekty nad bazą i pod bazą]
Niech $\mathbf 1$ będzie kategorią z jednym obiektem i jego identycznością, a $\Delta_B:\mathbf 1\to\mathcal C$ funktorem wskazującym $B$. Kategoria $(\id_{\mathcal C}\downarrow\Delta_B)$ jest kanonicznie izomorficzna z $\mathcal C/B$: obiektem jest strzałka $p:X\to B$, a warunek na morfizm redukuje się do $q\circ u=p$.

Kategoria $(\Delta_B\downarrow\id_{\mathcal C})$ opisuje obiekty pod $B$ i jest oznaczana $B/\mathcal C$. Jej obiektami są strzałki $p:B\to X$, a morfizm do $q:B\to Y$ jest strzałką $u:X\to Y$ z $u\circ p=q$. Obiektem początkowym jest $\id_B$, ponieważ jedyną strzałką spełniającą $u\circ\id_B=p$ jest $p$. W obu identyfikacjach usunięcie jedynego obiektu i morfizmu kategorii $\mathbf 1$ ma odwrotność polegającą na ich dołączeniu; obie operacje zachowują składanie i identyczności.
\end{przyklad}

Kategorie przecinkowe są przydatne przy opisie obiektów uniwersalnych i granic. W geometrii nie należy zakładać istnienia dowolnych iloczynów włóknistych rozmaitości; zob. przykład~\ref{cat:ex:manifold-pullback}.


